\chapter{Accounting methods for an operating company}
\label{ch:modules}

\section{A common core with distinct accounting books}

The common core governs evidence, time, money, events, journals and provenance.
Reporting standards govern recognition, measurement and presentation within
that core. The design preserves three levels of meaning: economic terms such
as a promise to pay cash; an accounting treatment under a particular standard;
and a company's reported label or extension. These levels are connected by
explicit rules and mappings, rather than collapsed into one universal chart
of accounts.

A rule package includes a standards register with jurisdiction, issuer class,
effective dates, transition provisions, permitted elections, paragraph
references, implementation version and tests. The package selector uses the
entity, reporting basis, period and adopted elections. It must distinguish a
standard's mandatory effective date from permitted early adoption. For example,
the IFRS Foundation's IFRS 18 overview states an effective date of 1 January
2027 with earlier application permitted \citep{ifrs18}. A 2026 historical
report cannot be assigned that presentation regime solely because the latest
installed package includes it.

IFRS, US GAAP and ASBE packages share economic contracts and events where
appropriate, but produce separate accounting books. A reconciliation between
books consists of explained adjustments with origins in differing treatments,
not a renaming of accounts. The first implementation can preserve all three
source bases even while only a selected subset has executable rules. Such a
source is represented as observed information with unsupported simulation
treatments, rather than being silently converted to the supported basis.

The official IFRS overviews for IFRS 9, IAS 32 and IFRS 17 establish the relevant
scope of financial-instrument measurement, presentation and insurance
contracts \citep{ifrs9,ias32,ifrs17}. They are orientation sources, not substitutes
for reviewed paragraph-level implementations. The retained official Ministry of Finance ASBE~21 page now supports selected
lease provisions. The official 2017 ASBE~22 text also supports the selected
credit-loss, interest and write-off distinctions worked through in
Chapter~\ref{ch:banking} \citep{mofasbe22}. Entity-specific issuance and amendment
chains and broader ASBE authority remain incomplete. Chapter~\ref{ch:insurance}
compares the dated ASBE~25 premium-allocation and incurred-claim provisions
with the retained AASB~17 text \citep{mofasbe25}. The Ministry's 2025
explanation supports the stated entity-class timetable, while Ping An's
reports supply an actual adoption and comparative-restatement example
(Section~\ref{sec:insurance-adoption}). Direct verification of the operative
issuance and implementation notices remains open, including the detailed
conditions of the later simplifications. Current ASC access is also unresolved.
Historical amendments and issuer policies inform proposed comparisons without
substituting for operative authority. These source limits are addressed
explicitly in the following dated-method discussion.

\section{Non-financial companies}

The initial industry package covers sales, collections, purchases, payments,
inventory movements, payroll, capital expenditure, depreciation and basic
financing. Revenue recognition is a separate rule from invoicing and cash
collection. A service contract may generate a performance event before an
invoice, creating a different asset from an unconditional receivable. Customer
advances create obligations before revenue. Returns, rebates and variable
consideration need separate assumptions and subsequent adjustments.

Inventory state includes quantity, cost, location, ownership and lot or cohort.
A costing method is selected within the reporting package's permissible
domain. An aggregate stock can support a simple model only if its approximation
is declared. A commodity-price shock can expose differences between old and
new inventory that an average value suppresses. Fixed-asset state likewise
requires acquisition dates, components, useful lives, residual values,
depreciation method, tax basis and impairment history when those affect the
requested forecast.

Industry extensions add economic mechanisms without bypassing the accounting
core. Construction requires performance and completion states; extractive
industries require reserves, depletion and restoration obligations; utilities
may require rate-regulation treatment; software businesses may require
contract liabilities and development-cost decisions. The module manifest
states what is represented, what is executable, and what remains aggregated.

\section{A matched inventory case across reporting profiles}

Standards packages need operation-level decisions. For interchangeable goods,
the inspected AASB~102 text incorporating IAS~2 permits FIFO or weighted
average in paragraph 25 and explains both in paragraph 27
\citep{aasb102}. The reproduced ASBE~1 text specifies FIFO, weighted average
or specific identification in Article 14; its official Ministry version could
not be retrieved in this review, so that reproduction has provisional source
status \citep{asbe1copy}. Caterpillar's company-authored 2024 report states
that it principally uses LIFO and supplies a FIFO comparison
\citep[Note 1.C]{caterpillar2024}. That is evidence of a disclosed US-GAAP
company policy; full ASC~330 paragraph authority remains an acquisition task.
It is not a rule that every US company must use LIFO.

Take the same purchases in three separate books: ten units at 5, followed by
ten at 7, then delivery of twelve units. FIFO charges $10(5)+2(7)=64$ to
cost of sales and leaves inventory 56. LIFO charges $10(7)+2(5)=80$ and
leaves inventory 40. Periodic weighted average uses cost $120/20=6$, charges
72 and leaves 48. The physical deliveries and supplier cash flows agree;
profit and carrying inventory differ by the selected costing rule. A concept
mapping cannot remove that difference. The FIFO/average selector in an
IAS~2 profile rejects LIFO as an unsupported election under the inspected
paragraph. The US company profile admits its evidenced LIFO policy while
retaining its narrower authority. An ASBE package remains source-provisional
until the attributed text and applicable amendments are officially verified.

Inventory impairment illustrates another required memory. Under inspected
AASB~102 paragraphs 28--34, an item costing 56 written down to net realisable
value 50 records expense 6. If the conditions reverse and net realisable
value becomes 54, reversal is 4, capped by original cost; the previous
write-down must therefore remain in state. The reproduced ASBE~1 Article 19
also addresses reversal within the earlier provision. This paragraph is not
used to assert the US treatment, whose primary authority was unavailable.
Source limitations constrain the treatment claim, while the separate-book
method and exact costing calculations remain fully specified.

\hypertarget{operating-events}{%
\section{Operating events}\label{method-corporate-methods-operating-events}}

The operating methods use the selected revenue, inventory and asset provisions of \citet{aasb15dated,aasb102,aasb116dated}. The lease and gross-tax case uses the dated lease, income-tax and cash-flow compilations \citep{aasb16dated,aasb112dated,aasb107dated}. Each rule still needs the entity, period, transaction facts and elections described below.

Every event uses the event-basis, authorization and commit contracts described in Section~\ref{sec:posting-permission}. The event inputs described here are required in addition to identity and date. Each object identifier resolves to its contract, context and admitted source or assumption. A purported cash collection does not establish a sale; a purchase order does not establish receipt.


Revenue, billing and collection are separate transitions. A delivery event names
the contract, performance obligation, quantity, allocated price and control
evidence. Recognition first releases a related contract liability; earned
consideration not yet unconditional creates a contract asset. Both credit
revenue. A missing control finding or unresolved price constraint leaves the
recognition unresolved.

An invoice or right-to-payment event names the contract, bill and unconditional
amount. It debits the receivable and credits an existing contract asset, with
any excess recorded as a contract liability. A collection names that receivable,
settlement and amount and debits cash while crediting the receivable. An excess
payment needs a separate customer-advance method; it cannot be credited to
revenue merely because cash arrived. These transitions retain cumulative
performance, billing and collection independently.

Purchasing starts when the company receives owned goods, not when it places an
order. The event records lot, quantity, unit cost and ownership evidence, debits
inventory and credits the payable. Consigned or unreceived goods cannot pass
that ownership condition. Consumption or sale identifies an available lot
allocation, quantity and admitted cost formula; it debits cost of goods sold
or an eligible production asset and credits inventory. Overselling does not
create stock. Supplier payment separately debits the payable and credits actual
cash. Insufficient cash blocks that settlement and can trigger a separately
modelled default; it never creates a hidden loan.

A tangible-asset event records component cost, residual value, useful life,
depreciation method and the date it becomes available for use. Initial debit to
PPE has a corresponding cash or payable credit; depreciation later debits
expense or an eligible asset and credits accumulated depreciation. Employee
service records earned pay, deductions and employer obligations for the service
period, separately from settlement. Unknown labour or tax law cannot become a
zero payable.

Closing derives revenue, expenses, OCI and owner movements from those events
using admitted statement mappings. It checks cash, inventory, financing and tax
roll-forwards rather than introducing a balancing entry. These event fields
form the method's required inputs; each identifier resolves to the same
contract, context and evidence used by its authorization.


AASB 15 paragraphs 9, 22, 31, 35--43, 47--59, 73--86 and 105--108 provide the selected revenue conditions; the retained compilation dated 31 December 2022 applies to annual periods beginning on or after 1 January 2023 and before 1 July 2026, covering the stated for-profit FY 2024 example. Variable consideration is estimated under a selected method and admitted only subject to the significant-reversal constraint, not recognized solely because its unconstrained expectation is positive. Principal versus agent, returns, warranties, contract modifications and distinct bundles require their own scoped decision records. If their facts are unresolved, retain alternative eligible methods; do not borrow the simple delivery rule by label.

The following schedules make the distinction between performance and billing explicit. For one performance obligation with price 100, cumulative revenue 60, unconditional bills 50 and collection 30, AR=20, CA=10, CL=0 and cash=30. Revenue 60 is represented by those three assets. Billing before performance changes intermediate balances but not the final result. For bills 80 instead, AR=50, CA=0, CL=20; net assets remain 60. Cumulative performed-service cost 40 then 70 posts 40 then 30, not 110. A downward revision follows the correction method and keeps the old result tied to its original view.

For this simple contract, let $R$ be cumulative recognised revenue, $B$
unconditional billings and $C$ collections, with $0\leq C\leq B$. The unpaid
bills determine accounts receivable (AR); performance in excess of billing is
a contract asset (CA), while billing in excess of performance is a contract
liability (CL). After each event,
\[
 \mathrm{AR}=B-C,\quad \mathrm{CA}=\max(R-B,0),\quad
 \mathrm{CL}=\max(B-R,0),\quad \text{cash received}=C.
\]
This recurrence applies per contract, with no impairment, variable consideration or unbilled prepayment; it is not permission to offset unrelated contracts. For billing 80 before delivery, debit AR 80/credit CL 80. Recognizing service 60 later debits CL 60 and credits revenue 60. Collecting 30 debits cash 30/credits AR 30, leaving AR 50 and CL 20, without a second receivable or an artificial contract asset. If service comes first, recognize CA 60/revenue 60, then bill AR 80 against CA 60 and CL 20. Both orders end at the same state. Cash received before an unconditional bill needs an explicit advance event that debits cash and credits CL and records the advance applied to later billing; it cannot pass through collect(AR).

Inventory uses identified cost for non-interchangeable items or the admitted cost formula for interchangeable items. AASB 102.9, 23, 25 and 28--34 support lower of cost/net realizable value, relevant cost formulas and write-down/reversal conditions. A method never chooses LIFO merely to fit a reported balance. In a second lot example, purchases are 10 units at 2 and 10 at 3, followed by sale of 12 units. FIFO gives cost 26 and stock 8 units/cost 24. Weighted average gives cost 30 and stock 20. Under the same given selling price 48, margin differs 22 versus 18. Both alternatives require their own permitted policy, consistent application and evidence; matching total cash does not remove the difference. NRV 16 for the remaining FIFO lot records expense 8, inventory allowance 8; recovery to NRV 22 reverses 6, leaving allowance 2. Recovery to 30 is capped at original cost 24. US LIFO remains a separately sourced question.

The retained AASB 116.7/15, 29--30 and 48--62 support cost and depreciation. Asset 100, residual 10, three-year even use gives 30 per full year. A year-two estimate change after first-year depreciation to residual 10 and three remaining years gives (70\ensuremath{-}10)/3=20 prospectively; an error in original cost is a correction, not an estimate change. Component replacements derecognize the replaced component's remaining amount before recognizing its replacement; no perpetual double asset. Borrowing-cost capitalization, impairment, provisions, employee benefits and share-based payment are specified as additional scoped methods below, not silently assigned zero in real-company cases.

\section{Apple: asset balances do not identify asset movements}
\label{sec:apple-operating-assets}

The preceding methods need receipts, lots, useful lives and settlement dates.
Apple's disclosures supply some of the balances those events must explain,
but they do not supply the events themselves. Its FIFO and straight-line
policies narrow the permitted cost calculations without identifying the actual
inventory lots or asset cohorts \citep[Note 1, printed p.~34]{apple2024}.
Comparing the notes with the cash-flow statement shows why filling those
missing inputs by subtracting balances can produce a wrong history.

\subsection{An inventory adjustment is not purchases or payments}

Between 28 September and 28 December 2024, components fall from 3,627 to 2,792
million dollars while finished goods rise from 3,659 to 4,119 million. Total
inventory falls from 7,286 to 6,911: the component changes, $-835$ and $+460$,
sum to $-375$ million \citep[Note 5, printed p.~9]{apple2025q1}. Those two
classes do not reveal quantities, unit costs or the transfers and sales that
produced the balances.

The quarter's cash-flow statement reports $+215$ million for inventories within
its reconciliation of net income to operating cash \citep[printed p.~5]{apple2025q1}.
That is an indirect reconciliation adjustment, not total cash paid to suppliers.
Even replacing it with minus the stock change would give $+375$, a different
quantity. The difference $215-375=-160$ million remains unexplained by these
selected disclosures. In fiscal 2024, inventory rose from 6,331 to 7,286, while
the cash-flow adjustment was $-1{,}046$; relative to the $-955$ stock-change
shortcut the difference is $-91$ million \citep[printed pp.~31 and 33]{apple2024}.
Neither residual is an identified currency effect, write-down or other journal
entry. Production, consumption, ownership, valuation and payment histories must
be established before they can explain the differences.

\subsection{Testing a proposed fixed-asset history}

The fixed-asset notes distinguish three measurements. The table shows gross
property, plant and equipment (PP\&E), its accumulated depreciation and the
remaining carrying amount, all in millions of dollars
\citep[Note 5, printed p.~39]{apple2024}
\citep[Note 5, printed p.~9]{apple2025q1}.

\begin{center}
\begin{tabularx}{\linewidth}{@{}Xrrr@{}}
\toprule
Amount & 30 Sep 2023 & 28 Sep 2024 & 28 Dec 2024\\\midrule
Gross PP\&E & 114,599 & 119,128 & 120,615\\
Accumulated depreciation & 70,884 & 73,448 & 74,546\\
Net PP\&E & 43,715 & 45,680 & 46,069\\\bottomrule
\end{tabularx}
\end{center}

Accumulated depreciation is shown here as a positive deduction: gross less
that deduction equals net. During fiscal 2024, gross cost increases 4,529,
accumulated depreciation increases 2,564 and net PP\&E increases 1,965.
Yet the note reports PP\&E depreciation expense of 8.2 billion, and the
cash-flow statement reports capital payments of 9,447 million. Its broader
depreciation-and-amortisation adjustment is 11,445 million
\citep[printed pp.~33 and 39]{apple2024}. These amounts differ in meaning;
the 8.2 billion also has coarser displayed precision. The change in accumulated
depreciation cannot be substituted for the year's depreciation charge.

Consider a deliberately limited asset book to test a tempting reconstruction.
Let $G_0,G_1$ be opening and closing gross cost and $B_0,B_1$ accumulated
depreciation, recorded as positive amounts. Let $A$ be gross additions, $R$
the original cost of retired assets, $d$ the depreciation charge and $S$ the
accumulated depreciation removed with those assets. Assume there are no other
movements and all depreciation is expensed. For this illustration, use the
displayed amounts as exact, including $d=8{,}200$ million. The two balances
must then satisfy
\begin{equation}
 G_1=G_0+A-R,\qquad B_1=B_0+d-S,\qquad 0\leq S\leq R.
 \label{eq:apple-asset-movement}
\end{equation}
The inequality requires a nonnegative carrying amount for the retired assets
in this cost model. If additions are equated with the 9,447 of cash payments,
the gross balance requires $R=9{,}447-4{,}529=4{,}918$. The depreciation
balance requires $S=8{,}200-2{,}564=5{,}636$. The proposed book removes more
accumulated depreciation than original cost, by 718 million. It is infeasible
under these combined assumptions. That finding rejects the proposed history;
it does not establish an error in Apple's filing or identify which assumption
must change.

The result is not sensitive to small rounding differences under a specified
rounding model. If the million-dollar balances and payment are rounded to the
nearest million, and 8.2 billion to the nearest 0.1 billion, the excess
$d-9{,}447+(45{,}680-43{,}715)$ lies in the closed outer envelope
$[666.5,769.5]$ million. Its lower endpoint is
$8{,}150-9{,}447.5+45{,}679.5-43{,}715.5=666.5$.
This is a conditional precision check, not a verified issuer rounding rule;
the exact tie convention remains unknown.

Allowing additions to be partly unpaid provides two simple repairs. Retain
the same selected opening and closing balances, 8,200 depreciation and 9,447
cash payments. Assume retirement at the start of the year with zero proceeds,
depreciation of surviving assets, and new acquisitions at year end. There are
no opening capital payables, taxes or other asset movements. The figures below
are constructed histories, not estimates of actual Apple transactions.

\begin{center}
\begin{tabularx}{\linewidth}{@{}Xrr@{}}
\toprule
Annual movement (USD million) & History 1 & History 2\\\midrule
Gross additions & 10,529 & 11,529\\
Gross cost retired & 6,000 & 7,000\\
Accumulated depreciation removed & 5,636 & 5,636\\
Retired carrying amount, expensed & 364 & 1,364\\
Cash paid for additions & 9,447 & 9,447\\
Additions remaining payable & 1,082 & 2,082\\\bottomrule
\end{tabularx}
\end{center}

Both histories satisfy equation~\eqref{eq:apple-asset-movement}; their gross
additions can also be allocated to match the three reported asset classes.
For History 1, acquisition debits PP\&E cost 10,529 and credits cash 9,447 and
a capital payable 1,082. Depreciation debits expense and credits accumulated
depreciation by 8,200. Retirement debits accumulated depreciation 5,636 and
loss 364, while crediting gross PP\&E cost 6,000. Net PP\&E rises 1,965 and
cash falls 9,447, so total assets fall 7,482. Liabilities rise 1,082 and equity
falls 8,564 through the two expenses: $1{,}082-8{,}564=-7{,}482$.
History 2 has the same asset change but liabilities rise 2,082 and equity
falls 9,564. These are coherent subsystem consequences, not a fit to all of
Apple's statements or an exhaustive family of real asset histories. Neither
history receives a probability or preference from this arithmetic.

The following quarter supplies a new snapshot, not the missing annual movements.
Its gross, accumulated and net changes are 1,487, 1,098 and 389 million,
whereas capital payments are 2,940 and broad depreciation/amortisation is
3,080 \citep[printed pp.~5 and 9]{apple2025q1}. The inspected interim note does
not separately identify PP\&E depreciation expense. The lease note adds a
related distinction: finance-lease right-of-use assets are included in PP\&E,
while operating-lease right-of-use assets are in other noncurrent assets.
The 1.0 billion of noncash right-of-use acquisitions covers the disclosed lease
population \citep[Note 8, printed p.~42]{apple2024}; adding it all to PP\&E
or to cash capital expenditure would erase that distinction.

The method therefore preserves reported balances, cash payments, expense and
unresolved movement components separately. A future cost or funding query
still needs lot and asset ages, useful lives, residual values, impairment and
settlement facts. The real filing has supplied constraints and ruled out one
simple history, while leaving several different explanations open.


\hypertarget{lease-and-tax-integration}{%
\section{Lease and tax integration}\label{method-corporate-methods-lease-and-tax-integration}}

A lease connects asset use, a financing obligation and income tax. The following schedule extends the operating-company example by separating these three measurements. The retained lease and tax dossiers supply the dated provisions, modification case and assumptions. The official lease compilation is dated December 2023 despite the older filename. The right-of-use (ROU) asset represents the admitted use of the leased asset. Its cost, the lease liability, accumulated depreciation, remaining payments, rate version and option judgments are separate state fields. Tax retains carrying amounts, legal tax bases, reversal dates, enacted/substantively-enacted rate evidence, recovery support, and original profit/OCI/equity allocation. The deferred tax asset (DTA) and deferred tax liability (DTL) are retained gross before any permitted offset presentation. A liability schedule is not a tax-base schedule.

For the two annual 121 payments at 10\%, initial liability and ROU are 210. Year-one interest is 21, liability after payment 110, depreciation 105 and ROU 105. Year two interest is 11, and both balances finish at zero. A same-asset modification just after year one creates two future payments 144 at revised 20\%, PV 220, so liability increases 110 and ROU increases 110 to 215. A payment index reset has different rate rules; it cannot reuse that modification blindly.

To demonstrate connection to company statements, assume a synthetic entity with opening cash/equity 500, no other opening balances, and annual cash service revenue 400 less cash operating costs 100, all recognized in their years. There are no dividends, other assets, debt, losses or permanent differences. Add the base lease above. Assumed tax law deducts cash lease payments against taxable income, attributes those deductions to the lease liability, gives zero tax bases to ROU/liability, and taxes profits at enacted 25\%. Sufficient usable profits and all recognition conditions for DTA are assumed. Tax is paid the following year. No current-period tax payment is invented from tax expense.

At commencement debit DTA 52.5 and credit DTL 52.5 in addition to the lease recognition. At year-one close DTA is 27.5 and DTL 26.25. The DTA reduction 25 is deferred expense; DTL reduction 26.25 is deferred benefit, net benefit 1.25. Taxable profit is 300\ensuremath{-}121=179; current tax/payable 44.75. Pretax book profit is 300\ensuremath{-}105\ensuremath{-}21=174, total tax 43.5 and profit 130.5. At year two, deferred expense 1.25 combines with current tax 44.75 to give total tax 46 on pretax profit 184.

\begin{center}
\begin{tabular}{lrr}
\toprule Closing balance & 2024 & 2025\\\midrule
Cash & 679 & 813.25\\
ROU asset & 105 & 0\\
Deferred tax asset & 27.5 & 0\\
Lease liability & 110 & 0\\
Deferred tax liability & 26.25 & 0\\
Current tax payable & 44.75 & 44.75\\
Equity & 630.5 & 768.5\\\bottomrule
\end{tabular}
\end{center}


2024 assets 811.5 equal liabilities 181 plus equity 630.5; 2025 assets 813.25 equal liabilities 44.75 plus equity 768.5. For this case, elect financing classification for lease interest under the retained AASB 107 regime. The operating and financing cash totals below depend on that presentation election. Direct operating cash is 300 then 255.25 after the first tax payment; financing lease cash is \ensuremath{-}121 each year; no investing cash arises. Indirect operating cash in year one is profit 130.5 plus depreciation 105 plus interest 21 minus deferred benefit 1.25 plus unpaid current tax 44.75 = 300. Year two gives profit 138 +105 +11 +1.25 =255.25. The second tax payable settles next year. Zero OCI and no owner transactions mean the equity change equals cumulative profit. This case tests gross tax, cash timing and financial-statement coherence, not a real entity's tax law.

\section{Apple: recognising tax and settling the obligation}
\label{sec:apple-tax-settlement}

The lease example separates a tax charge from the payment that follows it.
Apple's fiscal 2024 disclosures show why that separation matters when several
tax jurisdictions, deferred balances and a large disputed obligation enter
the same accounts. Its annual income-tax expense was 29,749 million dollars;
net cash paid for income taxes was 26,102 million
\citep[printed pp.~29 and 33]{apple2024}. The 3,647 difference is not, by itself,
an increase in tax payable. Deferred tax and other balance movements must first
be distinguished from current tax and settlement.

\subsection{The expense and balance reconciliations answer different questions}

Note 7 divides the 2024 expense between current and deferred amounts for federal,
state and foreign taxes. In millions of dollars, the reported components are
as follows \citep[Note 7, printed p.~40]{apple2024}. Parentheses in the source
represent negative expense, or a tax benefit.

\begin{center}
\begin{tabularx}{\linewidth}{@{}Xrrr@{}}
\toprule
Tax category & Current & Deferred & Total\\\midrule
Federal & 5,571 & $-3{,}080$ & 2,491\\
State & 1,726 & $-298$ & 1,428\\
Foreign & 25,483 & 347 & 25,830\\
Sum & 32,780 & $-3{,}031$ & 29,749\\\bottomrule
\end{tabularx}
\end{center}

The current charge is therefore 32,780 and the deferred benefit 3,031.
The tax payable has its own scope. Note 6 reports current income taxes payable
of 26,601 and noncurrent income taxes payable of 9,254 at 28 September 2024,
against 8,819 and 15,457 a year earlier. Combining both maturity classes gives
35,855 against 24,276, an increase of 11,579
\citep[Note 6, printed p.~39]{apple2024}. Current expense less net cash taxes
is only $32{,}780-26{,}102=6{,}678$. Reclassifying an amount between current
and noncurrent cannot explain the difference in their combined total.

To expose the missing movements, let $T$ denote that combined tax payable and
$N$ the net deferred tax asset reported in Note 7. Let $C$ be current tax
expense, $D$ deferred tax expense (negative for a benefit), and $P$ net cash
taxes paid, all for the same annual period and in millions of dollars. Define
$U$ and $V$ as the residuals needed to connect the reported amounts:
\begin{equation}
 \Delta T=C-P+U,\qquad \Delta N=-D+V.
 \label{eq:apple-tax-bridges}
\end{equation}
These are diagnostic definitions, not assertions that every tax balance has
the same recognition or settlement scope. The first relation gives
$U=11{,}579-6{,}678=4{,}901$. Setting that residual to zero would produce the
wrong closing payable. Assigning it to currency translation, uncertain tax
positions or another journal would require the missing transaction and
classification evidence.

The deferred balance illustrates the second relation. At September 2024,
gross deferred tax assets of 34,873 less a valuation allowance of 8,866 leave
26,007. Deducting deferred tax liabilities of 6,805 gives net deferred tax
assets of 19,202, compared with 17,251 a year earlier
\citep[Note 7, printed p.~41]{apple2024}. The increase is 1,951, so
$V=1{,}951-3{,}031=-1{,}080$. The deferred benefit in profit does not explain
the whole net-balance movement. Its allocation among currency, equity, other
comprehensive income and other causes remains unresolved here.

Note 6's 19,499 deferred-tax-asset presentation is also different from Note 7's
19,202 net position. Their difference of 297 is an arithmetic consequence,
not a separately disclosed liability line or authority to offset jurisdictions.
The reported temporary differences---including research expenditure, deferred
revenue, leases and depreciation---need their own tax bases and recovery
conditions. The note reports indefinitely available Irish foreign-tax and
California R\&D credit carryforwards, but also a valuation allowance for those
credits and some other differences. An indefinite term does not establish that
an asset is recoverable \citep[Notes 6--7, printed pp.~39--41]{apple2024}.

\subsection{A recognised charge, restricted assets and a later payment}

Apple attributes its large 2024 State Aid charge to the September 2024 court
decision confirming the European Commission's recovery decision. This is the
issuer's account of the dispute and its accounting consequence. Note 7 describes
a 15.8 billion dollar payable to Ireland, partly offset by a 4.8 billion US
foreign tax credit and an 823 million decrease in unrecognised tax benefits
\citep[Note 7, printed pp.~39--40]{apple2024}.

The narrative calls the net charge 10.2 billion; the statutory-rate reconciliation
gives the more precise amount 10,246 million. Subtracting the displayed narrative
centres gives $15{,}800-4{,}800-823=10{,}177$, which is 69 million lower.
That is compatible with the coarser inputs: if the two billion-dollar amounts
are rounded to the nearest 0.1 billion and 823 to the nearest million, the
component calculation lies in the closed outer interval
$[10{,}076.5,10{,}277.5]$ million, containing 10,246. The 10.2 billion charge
is compatible with it as well. This is a conditional rounding check, with the
issuer's precise rule and tie convention unverified; the centre difference
does not establish a conflicting source assertion.

The tax charge does not make the settlement funds disappear at recognition.
At 28 September 2024 the escrow comprised 2.6 billion dollars of cash and cash
equivalents and 13.2 billion of current marketable securities, restricted from
general use \citep[Notes 4 and 7, printed pp.~37 and 39]{apple2024}.
The annual liquidity discussion separately identifies the obligation and the
funds held to settle it. The following quarter reports both a payment and the
remaining obligation \citep[printed p.~26]{apple2024}
\citep[printed p.~16]{apple2025q1}.

\begin{center}
\begin{tabularx}{\linewidth}{@{}Xrr@{}}
\toprule
State Aid obligation and settlement & EUR billion & USD billion\\\midrule
Obligation at 28 Sep 2024 & 14.2 & 15.8\\
Payment during the following quarter & 11.0 & 11.9\\
Obligation at 28 Dec 2024 & 3.3 & 3.4\\\bottomrule
\end{tabularx}
\end{center}

The three rows are different dated quantities. In dollars, a movement equation
with opening obligation 15.8, payment 11.9 and closing obligation 3.4 requires
an additional net movement of $3.4+11.9-15.8=-0.5$ billion. Under nearest
rounding to 0.1 billion, its outer interval is $[-650,-350]$ million, excluding
zero. The same calculation in euros gives $3.3+11.0-14.2=0.1$ billion, with
outer interval $[-50,250]$ million, which includes zero. Neither calculation
supplies a common exchange rate or an exact exchange-gain entry. The interim
tax discussion mentions foreign-currency revaluation effects related to the
decision, but does not allocate this 500 million dollar residual
\citep[printed p.~16]{apple2025q1}.

The remaining funds also retain their asset classifications. At December 2024,
restricted cash and cash equivalents were 126 million dollars and restricted
current marketable securities 3.3 billion \citep[Note 4, printed p.~8]{apple2025q1}.
Their displayed sum, 3,426 million, is compatible with the rounded 3.4 billion
obligation. It is not all cash available for ordinary operations. The opening
2.6 billion of restricted cash alone could not have funded the reported
11.9 billion payment; securities sales, maturities, transfers and valuations
need their own histories.

For an isolated settlement illustration, suppose 11,900 million of restricted
cash is available and the amount paid equals the liability's carrying value
at settlement. Debit tax payable and credit restricted cash by 11,900. Assets
and liabilities each fall by 11,900; equity and current-period tax expense do
not change in this posting. This assumes the required cash conversion has
already occurred and omits any associated gain, loss or fee. It explains
settlement of an existing obligation without pretending to reconstruct Apple's
funding and currency entries. The quarter separately reports tax expense of
6,254 and net cash taxes of 18,651 million
\citep[printed pp.~1 and 5]{apple2025q1}; neither should be replaced with the
Ireland payment or treated as another copy of the preceding year's charge.

\subsection{Uncertain tax positions do not supply a payment distribution}

The last part of Note 7 preserves another population: gross unrecognised tax
benefits. Its 2024 movement is $19{,}454+1{,}727-386+2{,}542-1{,}070-229
=22{,}038$ million. The increases concern prior- and current-year positions;
the decreases concern changed prior-year positions, settlements and expired
limitation periods. Interest and penalties are excluded from this rollforward
\citep[Note 7, printed p.~41]{apple2024}. The 1,070 settlement-related decrease
is a movement in that gross balance, not an independently identified total of
cash paid to tax authorities.

Of the rounded 22.0 billion gross amount, the issuer says that 10.8 billion,
if recognised, would affect its effective tax rate. It also describes a
reasonably possible decrease of approximately 5--13 billion over the next
twelve months, primarily concerning transfer pricing and deemed repatriation
tax, while noting uncertainty about examination timing. These are different
conditional statements. The range is neither a probability law nor a promised
cash payment or expense release. A forecast still needs jurisdictional outcomes,
recognition, interest, penalties and settlement dates. Retaining those unknowns
keeps the tax note connected to the company's funding needs without turning
management's dated assessment into an invented future transaction.


\hypertarget{returning-to-the-company-forecast}{%
\section{Returning to the company forecast}\label{method-corporate-methods-returning-to-the-company-forecast}}

The Northstar case remains the broader ordinary-company teaching baseline: collections, purchases/payments, inventory, payroll, PPE, financing, tax, statements, annual/interim and revision views are already specified. Its 2024 collection/purchase/payment solution is 245/155/145. The 2025 forecast in the Northstar reference fixture opens with cash 74.6 and equity 90.6, spends 20 on capital expenditure at period end, pays dividend 8 and uses a base floating rate 7\%. Its simplified tax is settled in the year earned. The running case uses fixed-rate share $\alpha$; here floating-rate share $x=1-\alpha$ describes the same uncertainty. For 2025 the same unknown floating fraction $x\in[0,1]$ determines every
statement amount:
\begin{align*}
 I(x)&=5+2x, & \text{pretax profit}(x)&=82-2x,\\
 \text{tax}(x)&=16.4-0.4x, & \text{profit}(x)&=65.6-1.6x,\\
 \text{operating cash}(x)&=48.6-1.6x, & \text{closing cash}(x)&=95.2-1.6x,\\
 \text{equity}(x)&=148.2-1.6x, & \text{assets}(x)&=288.2-1.6x.
\end{align*}
Receivables 55, inventory 40, net PPE 98, payables 40 and debt 100 complete the
balance sheet. Liabilities are 140 for every $x$, so the common cash/equity
slope preserves the accounting identity. There is no default distribution on
$x$. At a 9\% floating rate the all-floating history has interest 9, profit
62.4, cash 92, equity 145 and assets 285. The annual cash bounds still cannot
establish daily payment solvency without event-date cash flows.

The lease/tax example extends the interaction evidence without changing the original Northstar fixture or pretending that its teaching 20\% tax formula is an IAS 12 method. In a real-company package, each required method is admitted separately and unsupported operations remain visible. The extended obligations for pensions, impairment, provisions, valuation and share-based payment are in the additional corporate methods that follow; omission from a finite company case never proves the actual company has none.


\hypertarget{employee-benefits}{%
\section{Employee benefits}\label{method-corporate-extensions-employee-benefits}}

The additional methods below use the selected IAS 19, IAS 36, IAS 37, IFRS 2, IFRS 13 and IAS 23 provisions in the original 2023 EU incorporating regulation \citep{eu20231803}. This is historical mechanism support; later amendments and actual entity facts remain separate applicability inputs.

DefinedContributionState retains service periods, promised contribution formula, contributions due/paid and applicable legal terms. For earned service requiring contribution 10, debit employee expense 10/credit payable 10; payment settles it. Calling a plan defined contribution needs the limited-obligation facts; a guaranteed benefit or additional funding obligation cannot disappear by name.

The defined-benefit obligation (DBO) measures accrued employee benefits under the stated actuarial assumptions. DefinedBenefitState retains employee/cohort service, benefit formula, projected salary, retirement/turnover/mortality assumptions, discount-curve currency/term, obligation, plan-asset fair value, asset ceiling, contributions, benefits and remeasurements. IAS 19.67--74 requires projected unit credit and the stated service attribution; .83 supplies discount-rate requirements. Assumptions are separately admitted measurements, not constants provided by the accounting rule. A one-employee toy plan promises 100 for each completed service year, paid after year two, with certain completion, zero salary growth and supported 10\% discount. At year-one end accrued benefit 100 due in one year has PV 1000/11. In year two that opening obligation accretes 100/11 and the second earned benefit adds 100; obligation 200 then settles 200. Merely using current salary times years without timing would give the wrong year-one obligation.

A separate funded-plan roll-forward supplies an independent statement example: opening DBO 100, plan assets 80 and net liability 20; current service 12; interest 5 on DBO, interest 4 on assets at 5\%; contributions 10 at year end; benefits 8 paid from plan at year end; actuarial obligation loss 7; actual asset return 6. DBO closes 100+12+5\ensuremath{-}8+7=116, assets 80+6+10\ensuremath{-}8=88, net liability 28. Expense is 12+1=13, OCI loss is 7\ensuremath{-}(6\ensuremath{-}4)=5. Debit service expense 12/credit net pension liability 12; debit interest expense 1/credit liability 1; debit OCI 5/credit liability 5; debit liability 10/credit cash 10. End liability 20+13+5\ensuremath{-}10=28. Plan-paid benefits reduce both DBO and plan assets and do not cause a second employer cash payment. IAS 19.113, .120--127 supports these components; remeasurement OCI is not recycled to profit under 122. A surplus requires the asset ceiling and applicable minimum funding analysis, not recognition of any calculated negative liability.

Unknown actuarial inputs give conditional obligations or measurement\_unresolved. Curtailment, settlement, amendment and multi-employer/insured arrangements need their own fact conditions and measurement-date remeasurement. A buyout is not settlement if obligations remain. Sensitivities vary one assumption only when that conditional meaning is explicit; correlated risks require joint scenarios.

\hypertarget{impairment-and-provisions}{%
\section{Impairment and provisions}\label{method-corporate-extensions-impairment-and-provisions}}

ImpairmentState retains asset/CGU boundaries, carrying amounts, allocated goodwill, corporate assets, recoverable-amount inputs and previous impairments/depreciation. IAS 36.9--10 determines indicators and mandatory annual tests. Recoverable amount is the higher of value in use and fair value less disposal costs (.18); it is not the lower of two pessimistic scenarios. VIU needs the supported cash flows, forecast horizon/growth, discount and exclusion of inappropriate financing/tax or uncommitted restructuring effects under 33--57. Compute conditional values if those estimates are unknown; do not infer them from historical profitability.

For a cash-generating unit (CGU), assets excluding goodwill 100 plus goodwill 20, recoverable amount 90: loss 30 first removes goodwill 20 and then reduces other assets 10 subject to their floors under 104--105. Debit impairment 30/credit goodwill 20/credit asset allowance 10. If the asset floor is 95 the requested reduction to 90 is impermissible; apply the remaining allocation/recognition conditions, not a forced plug. For a standalone cost-model asset with carrying 80, recoverable 60, recognize 20 loss. At a later measurement before further depreciation, supported recovery to 90 permits reversal only 20 to the no-impairment 80 ceiling (.117--119). Goodwill impairment is not reversed (.124). Different dates require the hypothetical unimpaired depreciation path, not the old nominal cost as cap.

ProvisionState retains obligating past event, present legal/constructive duty, outflow assessment, estimate distribution, settlement dates, risk/discount basis, reimbursements and evidence. IAS 37.14 requires present obligation, probable outflow and reliable estimate; .36--47 governs best estimate/risk/discount. An intended future operating expenditure without that duty is not a provision. Assume an established obligation to pay 121 in two years, deterministic amount, appropriate 10\% rate and no double-counted risk: recognize expense/provision 100, accrete interest 10 then 11, settle 121. The passage-of-time increase is borrowing cost under 60. If the obligation is an eligible asset restoration cost, initial debit follows that asset's method rather than always expense; later changes need the applicable restoration interpretation. Contingent liabilities/disclosures remain visible when recognition conditions are not met; no generic 0.5 threshold may substitute for the relevant standard's required analysis.

\hypertarget{share-based-compensation}{%
\section{Share-based compensation}\label{method-corporate-extensions-share-based-compensation}}

AwardState distinguishes equity versus cash settlement, grant/vesting/service dates, granted units, grant-date fair value, service/nonmarket/market conditions, expected units vesting, cumulative expense and changes/cancellations. IFRS 2.7--11, 15--23 supports the employee equity-settled method. With 100 units, grant value 3, two-year service condition and all expected to vest, year-one cumulative expense 150 and equity 150. At vesting only 80 satisfy the service condition: cumulative 240, second-year expense 90. Debit employee expense and credit equity. Failure of a market condition is different under 21; if the employee satisfies other vesting conditions, it does not automatically reverse the grant-date measured service expense. Unknown conditions cannot be represented by a single forfeiture percentage without their type and evidence.

Cash-settled awards remeasure the liability under 30--33. Same 100 expected units, half service and current unit fair value 4 gives liability/expense 200 at year one; 80 actual units and value 5 at settlement gives total 400, second expense 200, then debit liability 400/credit cash 400. Reusing the equity grant value 3 would be wrong. Actual award valuation requires the specified model/market inputs, option rights, volatility, dividends and calibration; the fair values 3/4/5 are case inputs. Modifications and settlement-choice arrangements require their own conditions, with the original minimum expense and increment analyzed under the applicable provisions rather than overwriting the original grant.

\hypertarget{fair-value-and-capitalization-boundaries}{%
\section{Fair value and capitalization boundaries}\label{method-corporate-extensions-fair-value-and-capitalization-boundaries}}

A ValuationRequest names the unit of account, market/date, asset/liability characteristics, principal/accessible market, technique, input hierarchy, counterparty/own non-performance risk where relevant, calibration and sensitivity. IFRS 13.24, .42--43 and .67--80 guide the selected measurement. Quoted identical asset price 7 for 10 units gives 70 in a qualifying Level 1 case; a blockage discount based only on holding size cannot reduce it to 63 under 69/80. A present-value illustration 111 due next year at an independently supported 11\% rate gives 100; that equality does not establish the rate, market comparability or accounting classification. Historical EIR, an impairment discount, risk-neutral pricing and a physical expected return answer different questions.

Borrowing-cost capitalization requires a qualifying asset, eligible costs, commencement/suspension/cessation dates and the specific/general-borrowing calculation, before moving interest into cost. IAS 23.8--25 supplies these rules in the retained historical EU text: for specific borrowing, eligible actual costs less temporary investment income; for general borrowing, the applicable weighted rate and cost ceiling. Commencement requires expenditures, borrowing costs and necessary preparation activities; extended inactive periods can suspend it, and substantial completion ends it. Assume specific borrowing 100 outstanding during a full eligible construction year at 5\%, with an admitted expenditure schedule and temporary investment income 1 on temporarily unused funds during that period. Net capitalized cost is 4. Real admission still needs the dated standard, actual schedule and eligibility evidence. It is not a default 5\% add-on to all assets. Internally developed intangibles likewise need research/development separation and all applicable recognition criteria; a forecast of future sales does not establish an asset.

The extension interface returns the stated measured components, proposed journal, economic delta, disclosures and unresolved inputs under the event-basis, authorization and commit contracts. Calculation success does not permit another operation to fill a missing authority/fact with a default. Real-company acquisition/tax/employee/valuation evidence and the full exceptions remain explicit internal preparation, with proposed answers and competing treatments retained for any later authorized domain review.



\hypertarget{rule-selection}{%
\section{Rule selection}\label{method-standards-methods-rule-selection}}

The proposed selection procedure distinguishes operative authority from historical source support. The retained sources include official AASB compilations, the original EU incorporating regulation and official ASBE lease and income-tax texts \citep{aasb9dated,aasb15dated,aasb16dated,aasb112dated,eu20231803,mofasbe21,moflease2018,mofasbe18,mofinterpretation16}. Historical US lease additions are separately identified in the retained amendment source \citep{fasbasu201602}.

AuthorityRequirement contains issuer/framework, jurisdiction, entity class, annual-period/application interval, paragraph and local source Ref, incorporating instrument, amendment chain, transition, elections and unresolved dependencies. RuleSelection binds those requirements to the accounting event and evidence view. The algorithm evaluates chronological applicability first, then scope and facts, then available elections/transition, then measurement/presentation. If two supported rules can apply, retain the explicit alternative and required election; if one is contradicted by facts, mark ineligible. Missing current authority gives unresolved, never eligibility through a familiar rule name.

An unchanged historical report remains an assertion under its own basis and vintage. A harmonized comparison creates an adjustment book and reconciliation, with changed recognition/measurement explained. Label mapping alone cannot convert US GAAP to IFRS or ASBE. The four clocks and source version remain part of each rule dependency; later knowledge cannot leak into an original as-of view. Transition entries reconcile opening assets, liabilities, tax and equity, rather than changing old source bytes or plugging an unexplained amount.

\hypertarget{ifrs-and-incorporating-jurisdictions}{%
\section{IFRS and incorporating jurisdictions}\label{method-standards-methods-ifrs-and-incorporating-jurisdictions}}

The retained official AASB texts provide explicit IAS/IFRS incorporation for their stated entity/version scope. AASB 9/15 compilation dates are 31 December 2022 with the source covers' application window; AASB 16 compilation 8 is 31 December 2023, AASB 112 compilation 7 is 29 September 2023. Their filenames are not date evidence. AASB 102's older retained edition is selected historical support for inventory, not current 2024 applicability certification. The retained EU 2023/1803 original provides official EU-adopted IFRS 3/10/IAS 19/21/36/37/IFRS 2/13 passages inspected in this method investigation; its historical adoption scope and later amendments remain separate. The retained amendment investigation supplies additional source evidence, not a blanket acceptance of all rules in these standards.

Initial real-method admission is therefore claim-specific: the selected revenue, lease/tax, financing, credit, group/FX and insurance cases must resolve each named applicability prerequisite. For synthetic demonstrations, declare the chosen dated provisions and assumptions and use the synthetic book mode. No assertion of whole-IFRS compliance follows from those entries.

The base lease in the saved dossiers (two annual 121 payments, 10\% rate, no qualifying exemptions or extra costs) has liability 210 initially, 110 after year one, 0 after year two. IFRS-equivalent cost-model ROU is 210/105/0; expenses 126 then 116. Income tax needs its own legal deductions and gross temporary differences, as the connected corporate case demonstrates. Revenue timing, inventory cost formula, financial instrument classification, control and contract boundaries are required rule predicates, not accounting options chosen after seeing which outcome fits best.

\hypertarget{us-gaap-proposed-methods-and-authority-limitation}{%
\section{US GAAP proposed methods and authority limitation}\label{method-standards-methods-us-gaap-proposed-methods-and-authority-limitation}}

The saved official FASB ASU 2016-02 Section A and later 2020-05/2021-09 amendments contain inspected historical ASC 842 additions. Their covers explicitly say an ASU is not authoritative. Current Codification access returned 403 again in this source investigation; the proposed methods cannot be marked current-authority-admitted. ASU amendments, historical FAS 109, issuer accounting policy and a model's memory are all insufficient substitutes for the operative ASC text and full applicable amendment chain. 

The same lease facts include ten-year remaining equipment life, fair value 1000, alternative use to lessor, no ownership transfer, no purchase option, two-year term and PV 210. Under the inspected historical 842-10-25-2/25-3 provisions, the stated facts support the proposed operating classification. Historical 842-20-25-6 and 35-3 yield single lease expense 121 per year, liability 210/110/0 and ROU 210/ 110/0. First year debit lease expense 121/credit liability 21/credit ROU 100; payment debits liability 121/credits cash 121. Next expense 121 decomposes into liability interest 11 and ROU reduction 110. The finance comparison uses the separate depreciation 105 plus interest 21/11 path. Changing classification facts is not a free election. The 75\%/90\% heuristics in the historical guidance cannot become universal requirements. This complete numerical contrast is useful despite the remaining current-authority limitation.

Proposed ASC 740 temporary-difference treatment retains tax base, reversal dates, enacted rates, valuation allowance and source-specific allocation/presentation. The saved historical FAS 109 and ASU 2019-12 provide inspected enactment support, not a completed current ASC 740 implementation. Under the matched PPE case, book 800/tax 600 at assumed 25\% yields DTL 50 and current tax 75 on taxable 300. A future 30\% rate that is merely substantively enacted but not legally enacted cannot be imported from the AASB 112 rule into this proposed US treatment; the declared 2025 closing temporary difference is 300, giving DTL 75 at the enacted 25\% rate versus DTL 90 at 30\%. This deferred-tax comparison is separate from current tax 75 on taxable profit 300. Actual legal enactment remains a separate jurisdictional fact.

The proposed CECL method for in-scope amortised-cost assets estimates lifetime expected cash shortfalls from origination using past/current/forecast evidence, with a documented reversion method beyond a reasonable/supportable forecast horizon and applicable instrument exceptions. It retains allowance roll-forwards, write-offs, recoveries, purchased-credit-deteriorated handling and off-balance commitments separately. It must not copy IFRS 9's 12-month/SICR staging to the same record. A candidate contrast: lifetime discounted loss 6 with 2 attributable to default events in the next year gives IFRS general Stage 1 candidate 2 versus CECL lifetime candidate 6, conditional on scope and method support. No current ASC 326 text was retrieved, and the attempted official amendment path returned 403; this contrast remains an explicitly unadmitted proposed reading, not a source-verified US requirement or posting example. Full current ASC 326, insurance 944, consolidation 810, FX 830 and relevant later amendments are specific internal source tasks before corresponding real admission.

\hypertarget{asbe-official-lease-support-and-remaining-tax-sources}{%
\section{ASBE lease adoption and deferred tax}\label{method-standards-methods-asbe-official-lease-support-and-remaining-tax-sources}}

For an ASBE lease, selecting the lease standard and selecting the related income-tax treatment require different application dates. Article 68 of the revised ASBE 21 states a general date of 1 January 2019, but the Ministry of Finance's 2018 issuance notice specifies the enterprise groups. Paragraph I requires application from that date for enterprises listed both domestically and overseas, and for overseas-listed enterprises preparing financial statements under IFRS or ASBE. Other enterprises applying ASBE must apply the revision from 1 January 2021 \citep[issuance notice, para.\ I]{moflease2018}.

Paragraph II permits earlier application for enterprises whose parent or subsidiary is listed overseas and prepares its overseas financial statements under IFRS or ASBE. Their early date cannot precede simultaneous application of the specified 2017 financial-instruments and revenue standards. Paragraph III discontinues the 2006 lease standard and its application guidance when the enterprise applies the revision. Thus a 2020 reporting date alone does not establish use of the revised lease standard for an enterprise in the second group: its qualifying early adoption would need evidence. The retained official Chinese texts include the 2018 notice and attachment and the Ministry's later ASBE 21 page, published on 28 October 2019 \citep{moflease2018,mofasbe21}. They replace the earlier secondary reproduction as support for the inspected wording.

Articles 4--8 require control of identified use, economic benefits and direction; article 6 includes both conditions for substantive substitution and the express presumption when difficult to determine. Articles 14--18 specify commencement, ROU cost and PV, using implicit rate or incremental borrowing rate if the former cannot be determined. Articles 20--23 separate cost-model depreciation and periodic interest; articles 25--29 distinguish reassessment/index changes/modifications and their discount-rate rules. Articles 30--33 define the exemptions; a two-year term is not short-term. Articles 60--67 specify transition choices and disclosures; the chosen transition cannot be inferred from a present closing liability \citep[attached standard, arts.\ 4--8, 14--33, 60--67]{moflease2018}.

For the declared base lease, these inspected recognition/measurement provisions support proposed ROU 210/105/0 and liability 210/110/0 with expense 126/116, under the same stated facts, depreciation assumption and fixed rule version. A same-asset modification is tested against 28--29, rather than labeled a separate lease without increased scope and commensurate consideration. A 2019 first-adopter using article 61(2) records cumulative opening changes without adjusting comparatives, while 61(1) is retrospective; comparison datasets must retain that distinction and the entity's actual election.

The tax question is why equal opening lease balances can still require separate deferred-tax assets and liabilities. The Ministry's original ASBE 18 defines tax bases and temporary differences in articles 4--9 and contains initial-recognition exceptions in articles 11 and 13. Article 13 also limits a deferred-tax asset to the probable taxable profit available to use the deduction. Articles 16--23 govern measurement, recovery or settlement assumptions, subsequent review, allocation and noncurrent classification \citep{mofasbe18}. The original text must be read with later interpretations.

Interpretation 16 removes the initial-recognition exemption for a single transaction that is not a business combination, initially affects neither accounting profit nor taxable income or loss, and creates equal taxable and deductible temporary differences. Section I(a) expressly includes qualifying initial lease liabilities and right-of-use assets; the corresponding deferred-tax liabilities and assets are recognized separately under ASBE 18. Equal carrying amounts alone do not establish equal temporary differences, and the asset's recoverability condition still applies. Section IV makes this treatment effective from 1 January 2023; Section I(b) permits early application from the year of issuance, 2022, with disclosure \citep[secs.\ I(a)--(b), IV]{mofinterpretation16}. This tax date is distinct from the enterprise's lease-standard adoption date. The operative text now supports the treatment previously illustrated by the retained official 2023 tax application case.

For the synthetic two-year lease, retain the assumed payments of 121, 10\% discount rate and straight-line depreciation. All amounts use the same illustrative monetary unit. Assume that tax deductions are allocated to the liability, both tax bases are zero, the applicable tax rate is 25\%, and sufficient taxable profit satisfies the deferred-tax asset recognition condition. Also assume that the transaction meets Section I(a)'s conditions and falls within its application period. At commencement the ROU asset and liability are each 210. The taxable ROU difference gives a deferred-tax liability (DTL) of 52.5; the deductible liability difference gives a deferred-tax asset (DTA) of 52.5. The entry debits DTA and credits DTL by those amounts. At the first annual close, liability 110 and ROU 105 give DTA 27.5 and DTL 26.25. Reducing the DTA by 25 creates deferred-tax expense; reducing the DTL by 26.25 creates a benefit, leaving net benefit 1.25. Both balances reverse by expiry. Applying the old exemption to omit both opening deferred balances is wrong for these qualifying facts and dates: it conceals the required gross balances and the later tax effect.

Transition also needs more than an opening journal entry. Section I(b) requires adjustments for qualifying transactions between the earliest comparative opening and application. For qualifying existing balances at that opening, cumulative effects adjust opening retained earnings and the relevant financial-statement items. Consider a separate transition example with cash 500, ROU 105, lease liability 110 and equity 495, with no deferred tax yet recognized. Under the same assumed tax bases, rate and recoverability, recognition debits DTA 27.5, credits DTL 26.25 and credits opening retained earnings 1.25. The resulting balance sheet is
\[
\underbrace{500+105+27.5}_{\text{assets }632.5}
=\underbrace{110+26.25}_{\text{liabilities }136.25}
+\underbrace{495+1.25}_{\text{equity }496.25}.
\]
The equity adjustment records the cumulative effect at the comparative opening. If the remaining lease year falls within the comparative year, the DTA of 27.5 and DTL of 26.25 reverse during that year, giving comparative deferred-tax expense 1.25. The opening adjustment alone therefore does not complete the transition. Neither ASBE 21's article 61(2) lease-adoption election nor Interpretation 16 Section III's no-comparative-adjustment rule for its separate share-payment topic may be substituted for Section I(b)'s tax transition.

These are gross recognition and measurement amounts. Separate recognition and noncurrent classification do not determine whether an actual balance sheet may offset DTA and DTL. That decision needs the applicable presentation provisions and entity/tax-authority facts. Actual tax law, tax bases and rates, recoverable profit, the enterprise's adoption and elections, and later amendments remain to be established before applying the example to an issuer. Full ASBE 22 credit, consolidated statements, group combinations and insurance revisions also retain specific source dependencies; similar amounts do not establish full IFRS equivalence.

\hypertarget{matched-case-outcomes-and-decision}{%
\section{Matched-case outcomes and decision}\label{method-standards-methods-matched-case-outcomes-and-decision}}

\begin{longtable}[]{@{}ll@{}}
\toprule
\begin{minipage}[b]{0.47\columnwidth}\raggedright
Input changed or withheld\strut
\end{minipage} & \begin{minipage}[b]{0.47\columnwidth}\raggedright
Required consequence\strut
\end{minipage}\tabularnewline
\midrule
\endhead
\begin{minipage}[t]{0.47\columnwidth}\raggedright
Framework changes with same lease facts\strut
\end{minipage} & \begin{minipage}[t]{0.47\columnwidth}\raggedright
Preserve different expense/ROU paths and source strength, not just rename accounts\strut
\end{minipage}\tabularnewline
\begin{minipage}[t]{0.47\columnwidth}\raggedright
Method source present but target entity's adoption unknown\strut
\end{minipage} & \begin{minipage}[t]{0.47\columnwidth}\raggedright
Unresolved applicability, no real posting authorization\strut
\end{minipage}\tabularnewline
\begin{minipage}[t]{0.47\columnwidth}\raggedright
ROU/liability initially equal\strut
\end{minipage} & \begin{minipage}[t]{0.47\columnwidth}\raggedright
Still evaluate gross tax differences and later divergent balances\strut
\end{minipage}\tabularnewline
\begin{minipage}[t]{0.47\columnwidth}\raggedright
Historical rule supports method but a later reporting year is requested\strut
\end{minipage} & \begin{minipage}[t]{0.47\columnwidth}\raggedright
Recheck operative amendments/transition; no automatic extension\strut
\end{minipage}\tabularnewline
\begin{minipage}[t]{0.47\columnwidth}\raggedright
Issuer policy discloses LIFO, regulator text unavailable\strut
\end{minipage} & \begin{minipage}[t]{0.47\columnwidth}\raggedright
Retain issuer assertion; do not upgrade it into universal US authority\strut
\end{minipage}\tabularnewline
\begin{minipage}[t]{0.47\columnwidth}\raggedright
Foreign basis label replaced by local label\strut
\end{minipage} & \begin{minipage}[t]{0.47\columnwidth}\raggedright
Reject as a standards conversion without recognition/measurement bridge\strut
\end{minipage}\tabularnewline
\bottomrule
\end{longtable}

These decisions establish a common procedure and specific proposed treatments. The recovered ASBE lease notice and tax interpretation support the stated application schedules and conditional example. Broader authority coverage remains incomplete: operative US sources, ASBE sector methods, actual tax law and presentation conditions, and the applicable amendment and entity-election evidence still require internal investigation.

